\documentclass[10pt,english]{article}

\usepackage{amsfonts,url,tikz,amsmath}

% Paper setup
\evensidemargin=0in
\oddsidemargin=0in
\textwidth=6.25in
\topmargin=-0.5in
\headheight=0.0in
\headsep=0.5in
\textheight=9.0in
\footskip=0.5in

\input{institution}

\newcommand{\Interior}[1]{\ensuremath{{#1}^{\circ}}}
\newcommand{\Closure}[1]{\ensuremath{\overline{#1}}}
\newcommand{\Complement}[1]{\ensuremath{{#1}^{c}}}

\newcommand{\Expect}{\ensuremath{\mathrm{E}}}
\newcommand{\vecnot}{\underline}
\newcommand{\RealNumbers}{\ensuremath{\mathbb{R}}}
\newcommand{\RationalNumbers}{\mathbb{Q}}
\newcommand{\ComplexNumbers}{\mathbb{C}}
\newcommand{\Real}{\mathrm{Re}}
\newcommand{\Span}{\mathrm{span}}
\newcommand{\Rank}{\mathrm{rank}}
\newcommand{\Nullity}{\mathrm{nullity}}
\newcommand{\Trace}{\mathrm{tr}}
\newcommand{\Diag}{\mathrm{diag}}
\newcommand{\dd}{\mathrm{d}}
\DeclareMathOperator*{\esssup}{ess\,sup}

% Use < , > inner product
\newcommand{\inner}[2]{{\left\langle #1 \mskip2mu , #2 \right\rangle}}
\newcommand{\tinner}[2]{{\langle #1 \mskip1mu , #2 \rangle}}

% Use < | > inner product
%\newcommand{\inner}[2]{{\left\langle #1 \mskip2mu \middle| \mskip2mu #2 \right\rangle}}
%\newcommand{\tinner}[2]{{\langle #1 \mskip1mu | \mskip1mu  #2 \rangle}}



\begin{document}

\title{ECE 586: Vector Space Methods \\ Lecture 9: Frequently Asked Questions}
\author{Henry D. Pfister \\ \univ}
\date{October 18, 2026}

\maketitle

\section{Fields and Matrices}

\paragraph{Why do we introduce fields when most engineering problems use real or complex numbers?}
A field records the arithmetic rules needed for scalar calculations, so the same argument can work over $\mathbb{Q}$, $\mathbb{R}$, $\mathbb{C}$, or a finite field.
Slides 1--3 separate these rules from any particular interpretation of the scalars.
For example, Gaussian elimination uses field arithmetic and therefore applies to linear systems over finite fields as well.

\paragraph{Why are the integers not a field? Why does arithmetic modulo a prime work?}
A field requires every nonzero element to have a multiplicative inverse.
The integer $2$ has no integer inverse, so $\mathbb{Z}$ is not a field.
Modulo a prime $p$, each nonzero residue has an inverse.
For a composite modulus this fails: modulo $6$, $2\cdot3=0$, although both factors are nonzero.

\paragraph{Does multiplication by zero have an inverse?}
No. The inverse axiom on slide 2, ``Fields,'' excludes zero.
Distributivity implies $0s=0$ for every scalar $s$, so $0s$ cannot equal the field identity $1\neq0$.
Division by zero is therefore undefined.

\paragraph{Why is the matrix for $m$ equations in $n$ unknowns an $m\times n$ array?}
Each row contains the coefficients of one equation, and each column corresponds to one unknown.
Thus $A\in F^{m\times n}$ multiplies a column vector $x\in F^n$ to produce $Ax\in F^m$, as on slide 4, ``Matrices.''

\paragraph{How should I interpret matrix multiplication beyond the entry formula?}
On slide 5, ``Matrix Multiplication,'' each column of $AB$ is a linear combination of columns of $A$, with coefficients from the corresponding column of $B$.
Equivalently, $AB$ represents first applying $B$ and then $A$.
The shared inner dimension is required for these combinations to be defined; multiplication need not commute.

\section{Row Reduction}

\paragraph{What is the difference between row echelon form and reduced row echelon form?}
In row echelon form, zero rows are last and each successive nonzero row starts strictly farther to the right.
Reduced row echelon form also requires each leading entry to be one and every other entry in its column to be zero.
Slide 6 illustrates both. The RREF of a given matrix is unique, although elimination can take different sequences of steps.

\paragraph{Why must a row be multiplied by a nonzero scalar?}
Elementary row operations must be reversible to preserve an equivalent linear system.
Multiplication by a nonzero scalar can be undone by its inverse, but multiplication by zero destroys an equation.
For $Ax=b$, apply every row operation to the entire augmented matrix $[A\;b]$, not only to $A$.

\paragraph{How does the augmented matrix $[A\;I]$ keep track of elimination?}
The same row operations that turn $A$ into $R$ turn $I$ into an invertible matrix $P$, giving $[R\;P]$ with $R=PA$.
On slide 8, only the left block must reach RREF; continue no farther once it has done so.
Row reduction preserves the null space because $PAx=0$ if and only if $Ax=0$.

\paragraph{Why does $m<n$ guarantee a nonzero solution of $Ax=0$?}
There can be at most $m$ pivot columns, so with $n>m$ at least one variable is free.
Choose a free variable to be nonzero and solve backward for the pivot variables.
For example, $[1\;2]x=0$ has the nonzero solution $x=(-2,1)^T$.
This is the ``Foundation of Dimension'' lemma on slide 9.

\paragraph{Does having more unknowns than equations guarantee a solution of $Ax=b$?}
No. It guarantees a nonzero homogeneous solution, not consistency for an arbitrary right-hand side.
For example, two identical coefficient rows with different right-hand sides are inconsistent.
If a solution $x_0$ does exist, all solutions are $x_0+z$ with $z\in\mathcal{N}(A)$.

\end{document}
